\chapter{Entraîner}

\section{La commande minimale}

Deux informations suffisent : le modèle de départ et le corpus.

\begin{nkcode}[]{Windows — PowerShell}
.\Build\Bin\Release-Windows\NKQwen2Train\NKQwen2Train.exe `
  --gguf="$env:USERPROFILE\.ollama\models\blobs\sha256-VOTRE_BLOB" `
  --corpus="C:\mes-donnees\corpus.txt" `
  --out="C:\mes-donnees\mon-modele"
\end{nkcode}

\begin{nkcode}[]{Linux et macOS}
./Build/Bin/Release-Linux/NKQwen2Train/NKQwen2Train \
  --gguf=~/.ollama/models/blobs/sha256-VOTRE_BLOB \
  --corpus=~/mes-donnees/corpus.txt \
  --out=~/mes-donnees/mon-modele
\end{nkcode}

C'est la même commande partout : seuls les chemins changent.

\section{Toutes les options}

\begin{center}
\begin{longtable}{p{3.9cm}p{2.0cm}p{7.6cm}}
\toprule
\textbf{Option} & \textbf{Défaut} & \textbf{À quoi ça sert} \\
\midrule
\endhead
\texttt{-{}-gguf=} & — & Le modèle de départ. \textbf{Obligatoire.} \\
\texttt{-{}-corpus=} & — & Le fichier de textes. \textbf{Obligatoire.} \\
\texttt{-{}-out=} & \texttt{qwen2\_lora} & Préfixe des fichiers produits. \\
\texttt{-{}-lr=} & \texttt{1e-4} & Taux d'apprentissage : l'ampleur de chaque
   correction. \\
\texttt{-{}-steps=} & \texttt{0} & Nombre de pas sur ce lancement. \texttt{0}
   signifie « jusqu'au bout du corpus ». \\
\texttt{-{}-save-every=} & \texttt{200} & Sauvegarde tous les $n$ pas. \\
\texttt{-{}-rank=} & \texttt{8} & Taille des pièces LoRA. \\
\texttt{-{}-max-seq=} & \texttt{128} & Longueur maximale d'un exemple, en
   tokens. Les plus longs sont sautés. \\
\texttt{-{}-valid-every=} & \texttt{500} & Contrôle tous les $n$ pas.
   \texttt{0} pour ne jamais contrôler. \\
\texttt{-{}-fresh} & — & Ignorer les checkpoints et repartir de zéro. \\
\bottomrule
\end{longtable}
\end{center}

\begin{nkretenir}
Tout se règle par la ligne de commande. \textbf{Changer de corpus, de taux
d'apprentissage ou de durée ne demande jamais de recompiler.} Recompiler ne
redevient nécessaire que si l'on modifie le programme lui-même.
\end{nkretenir}

\section{Lire ce qui s'affiche}

\subsection{Au démarrage}

\begin{nkcode}[]{Sortie réelle, machine à 8 Go}
=== NKQwen2Train — entraînement LoRA reprenable ===

backend compute : Vulkan
  [NkQwen2LoraGpu] L=28 d=3584 ffn=18944 nH=28 nKV=4 hd=128 vocab=152064
    couche 1/28 chargée (0.7 s, 0.31 Go)
    couche 28/28 chargée (10.1 s, 3.83 Go)
    lm_head = output.weight (Q6_K, 152064 x 3584, 426 Mo)
    adaptateurs : 20185088 paramètres entraînables (77 Mo) + 231 Mo d'état Adam
modèle chargé en 13.4 s
corpus : 100017 paires
aucun checkpoint exploitable — départ de zéro
\end{nkcode}

Ligne par ligne :

\begin{itemize}
  \item \textbf{backend compute : Vulkan} — la vérification du chapitre~2. Si
        cette ligne dit autre chose, arrêtez-vous.
  \item \textbf{L=28 d=3584} — 28 couches, 3584 nombres pour représenter un
        token. C'est la carte d'identité du modèle.
  \item \textbf{3,83 Go} — la place prise sur votre carte graphique. Sur 8~Go,
        il reste de quoi travailler.
  \item \textbf{20185088 paramètres entraînables} — les fameux 0,27~\%.
  \item \textbf{231 Mo d'état Adam} — la mémoire de l'optimiseur ; on y revient
        au chapitre~5.
  \item \textbf{100017 paires} — combien d'exemples ont été lus. Comparez à ce
        que vous attendiez.
\end{itemize}

\subsection{Pendant l'entraînement}

\begin{nkcode}[]{Une ligne toutes les dix étapes}
pas 10 | perte 2.2903 (moy. 1.8772) | 8.46 s/pas | exemple 10/100017
\end{nkcode}

\begin{center}
\begin{tabular}{lp{9.6cm}}
\toprule
\texttt{pas 10} & Nombre total de pas depuis le tout début, reprises
                  comprises. \\
\texttt{perte 2.2903} & L'erreur sur \emph{cet} exemple. Elle saute beaucoup :
                  c'est normal. \\
\texttt{(moy. 1.8772)} & La moyenne. \textbf{C'est celle-ci qu'on surveille.} \\
\texttt{8.46 s/pas} & Le rythme. \\
\texttt{exemple 10/100017} & Où l'on en est dans le corpus. \\
\bottomrule
\end{tabular}
\end{center}

\begin{nkpiege}
Ne tirez aucune conclusion de la perte d'un exemple isolé. Elle dépend d'abord
de la difficulté de cet exemple-là. Une phrase courte et banale donne une perte
basse, une phrase technique une perte haute — sans que le modèle ait progressé
ou régressé entre les deux. \textbf{Seule la moyenne, sur des centaines de pas,
raconte quelque chose.}
\end{nkpiege}

\section{À quoi ressemble un apprentissage qui marche}

Voici des mesures réelles, prises sur 200 pas :

\begin{center}
\begin{tabular}{rrr}
\toprule
\textbf{Pas} & \textbf{Perte d'apprentissage} & \textbf{Perte de contrôle} \\
\midrule
0   & —      & 1,6613 \\
1   & 2,2370 & 1,6212 \\
25  & 1,2378 & 0,8662 \\
50  & 1,8171 & 0,7430 \\
100 & 1,2766 & 0,6655 \\
150 & 0,9647 & 0,6580 \\
200 & 0,8639 & \textbf{0,6292} \\
\bottomrule
\end{tabular}
\end{center}

Deux choses à observer.

\textbf{La colonne du milieu monte et descend.} Au pas~50 elle vaut 1,82, plus
qu'au pas~25. Ce n'est pas une régression : c'est un exemple plus difficile.

\textbf{La colonne de droite descend franchement}, de 1,66 à 0,63. Comme elle
est mesurée sur des exemples jamais vus, cette baisse est un vrai progrès et non
de la récitation.

\begin{nkretenir}
Un bon entraînement se reconnaît à ceci : la perte de contrôle descend
régulièrement. Si elle stagne pendant des centaines de pas, quelque chose ne va
pas — le plus souvent le corpus, parfois le taux d'apprentissage.
\end{nkretenir}

\section{Combien de temps cela va-t-il prendre ?}

La question que tout le monde pose. La réponse mesurée :

\begin{center}
\begin{tabular}{lr}
\toprule
Durée d'un pas (moyenne) & 9,15 s \\
Durée d'un pas (observé) & de 8,5 à 14,2 s \\
\midrule
1\,000 pas & $\approx$ 2 h 30 \\
10\,000 pas & $\approx$ 25 h \\
Une époque complète (100\,017 exemples) & \textbf{254 h $=$ 10,6 jours} \\
\bottomrule
\end{tabular}
\end{center}

Oui : dix jours et demi pour un seul passage sur le corpus entier.

Ce chiffre n'est pas un défaut à corriger, c'est la réalité d'un modèle de sept
milliards de paramètres sur une carte grand public. Trois façons de faire avec :

\begin{enumerate}
  \item \textbf{Laisser tourner.} C'est pour cela que la reprise existe
        (chapitre~5). La machine travaille la nuit, vous l'arrêtez quand vous
        en avez besoin, vous la relancez ensuite.
  \item \textbf{Réduire le corpus à ce qui compte.} 10\,000 exemples bien
        choisis valent mieux que 100\,000 exemples médiocres — et prennent un
        jour au lieu de dix.
  \item \textbf{Écourter les exemples.} Baisser \texttt{-{}-max-seq} accélère
        chaque pas, au prix des exemples longs qui seront sautés.
\end{enumerate}

\begin{nkpiege}
La durée d'un pas dépend de la \textbf{longueur} de l'exemple. Un corpus de
questions courtes ira nettement plus vite qu'un corpus de longs paragraphes. Les
10,6~jours sont une moyenne sur \emph{ce} corpus-là, pas une loi universelle :
mesurez sur le vôtre avec \texttt{-{}-steps=20}.
\end{nkpiege}

\section{Régler le taux d'apprentissage}

C'est l'ampleur de chaque correction. Trop petit, le modèle n'apprend rien ;
trop grand, il part dans le décor.

\begin{center}
\begin{tabular}{lp{9.4cm}}
\toprule
\texttt{1e-3} & Rapide et risqué. La perte peut exploser. \\
\texttt{1e-4} & \textbf{Le défaut.} Fonctionne dans la grande majorité des cas.
                Commencez par là. \\
\texttt{1e-5} & Prudent. Pour affiner un modèle déjà bien adapté. \\
\bottomrule
\end{tabular}
\end{center}

\begin{nkpiege}
Si votre perte grimpe au lieu de descendre, ou affiche \texttt{nan}, votre taux
est trop grand. Divisez-le par dix et relancez avec \texttt{-{}-fresh} : les
paramètres déjà abîmés ne se récupèrent pas.
\end{nkpiege}

\section{Régler le rang}

Le rang fixe la taille des pièces ajoutées.

\begin{center}
\begin{tabular}{rrp{7.4cm}}
\toprule
\textbf{Rang} & \textbf{Paramètres} & \textbf{Quand le choisir} \\
\midrule
4  & $\approx$ 10~M & Adapter le style, peu de mémoire disponible. \\
8  & 20~M & \textbf{Le défaut.} Le bon compromis. \\
16 & $\approx$ 40~M & Domaine riche, beaucoup d'exemples. \\
32 & $\approx$ 80~M & Rarement utile ; à réserver aux très gros corpus. \\
\bottomrule
\end{tabular}
\end{center}

\begin{nkpiege}
Le rang doit être \textbf{identique} entre l'entraînement et la reprise. Reprendre
un checkpoint de rang~8 en demandant un rang~16 est refusé — et c'est heureux :
les formes ne correspondraient pas, et le modèle rechargé serait faux.
\end{nkpiege}

\section{Ce que vous obtenez à la fin}

\begin{nkcode}[]{Fin d'un entraînement}
checkpoint final : mon-modele_ckpt0.nkla (231.0 Mo)
adaptateurs livrables : mon-modele.nkla (77.0 Mo, 20185088 paramètres)

6 pas en 76 s (12.73 s/pas) | perte moyenne 1.9923 | total cumulé : 6 pas
\end{nkcode}

Deux sortes de fichiers, à ne pas confondre :

\begin{center}
\begin{tabular}{p{4.3cm}rp{7.2cm}}
\toprule
\textbf{Fichier} & \textbf{Taille} & \textbf{À quoi il sert} \\
\midrule
\texttt{\_ckpt0.nkla}, \texttt{\_ckpt1.nkla} & 231 Mo & \textbf{Reprendre} un
   entraînement. Contient tout l'état. \\
\texttt{.nkla} & 77 Mo & \textbf{Utiliser} le modèle. C'est ce fichier-là que
   vous partagez. \\
\bottomrule
\end{tabular}
\end{center}

\begin{nkretenir}
Le fichier à donner, c'est le \texttt{.nkla} de 77~Mo. Les checkpoints ne
servent qu'à vous, sur votre machine, pour continuer là où vous vous êtes
arrêté.
\end{nkretenir}

\begin{nkexo}
Reprenez le corpus de vingt paires du chapitre~3 et lancez :
\texttt{-{}-steps=20 -{}-save-every=10 -{}-valid-every=0}. Comptez le temps
qu'un pas prend \emph{sur votre machine} et multipliez par le nombre d'exemples
de votre vrai corpus. Vous saurez à quoi vous engager avant de vous engager.
\end{nkexo}
